\section{Proofs for Section~\ref{sec:heights}}
\label{app:heights}

This appendix proves the results of Section~\ref{sec:heights}. We keep the
notation of that section: $z(X)$ is the vector of monomials of degree at most
two with the constant first, $N=\binom{n+2}{2}$, $w(X,v)=(v,X\otimes v)$, and
a certified quartic is a rational quartic with a rational positive definite
full Hessian Gram. For a polynomial $P$ of degree at most two,
$\operatorname{coef}(P)\in\R^N$ denotes its coefficient vector, so that
$P=\operatorname{coef}(P)^{\mathsf T}z$. For a rational $a/b$ in lowest terms
we use $\bits(a/b)=\lceil\log_2(|a|+1)\rceil+\lceil\log_2(b+1)\rceil$, which is
at least $\log_2\max\{1,|a|\}+\log_2b$.

\subsection{The Taylor Gram}
\label{app:heights-taylor}

\begin{proof}[Proof of Lemma~\ref{lem:taylor-sos}]
\emph{(a)} For fixed $X$ the polynomial $\varphi(t)=f(a+t(X-a))$ satisfies
$\varphi(1)=\varphi(0)+\varphi'(0)+\int_0^1(1-t)\varphi''(t)\,dt$. With
$d=X-a$ this reads
\[
 f(X)=f(a)+b^{\mathsf T}d+\int_0^1(1-t)\,d^{\mathsf T}\nabla^2f(a+td)\,d\,dt.
\]
By \eqref{eq:heights-hessian-gram} with the point $a+td$ and the direction
$d$, the integrand is $(1-t)\,w(a+td,d)^{\mathsf T}A\,w(a+td,d)$, and
$w(a+td,d)=(d,(a+td)\otimes d)=U+tV$. Since
$\int_0^1(1-t)\,dt=\frac12$, $\int_0^1(1-t)t\,dt=\frac16$, and
$\int_0^1(1-t)t^2\,dt=\frac1{12}$, the integral equals
\[
 \tfrac12U^{\mathsf T}AU+\tfrac13U^{\mathsf T}AV+\tfrac1{12}V^{\mathsf T}AV
 =\tfrac12(U+V/3)^{\mathsf T}A(U+V/3)+\tfrac1{36}V^{\mathsf T}AV,
\]
where symmetry of $A$ was used. Substituting $U=C_U(a)z(X)$ and
$V=C_V(a)z(X)$ gives
$z(X)^{\mathsf T}\mathcal T_a(A)z(X)$. The same computation shows, for every
symmetric $B$ of order $n+n^2$,
\begin{equation}
 z(X)^{\mathsf T}\mathcal T_a(B)z(X)
 =\int_0^1(1-t)\,(U+tV)^{\mathsf T}B\,(U+tV)\,dt.
 \label{eq:heights-taylor-form}
\end{equation}

\emph{(b)} The matrix $\mathcal T_a(A)$ is a sum of congruences of $A$ with
positive scalar factors, so it is positive semidefinite when $A$ is. Suppose
$A\succ0$. Then $\xi^{\mathsf T}\mathcal T_a(A)\xi=0$ holds exactly when
$(C_U(a)+\frac13C_V(a))\xi=0$ and $C_V(a)\xi=0$, that is, when
$C_U(a)\xi=0$ and $C_V(a)\xi=0$. The rows of $C_U(a)$ and $C_V(a)$ are the
coefficient vectors of the polynomials $d_i$, $a_id_j$, and $d_id_j$ for
$i,j\in\{1,\ldots,n\}$. Each of them vanishes at $a$, and the $d_i$ together
with the $d_id_j$, $i\le j$, form a basis of the hyperplane
$\mathcal H_a=\{P:\deg P\le2,\ P(a)=0\}$, since they are the nonconstant
monomials in the translated variables. Because
$P(a)=\operatorname{coef}(P)^{\mathsf T}z(a)$, the coefficient vectors of
$\mathcal H_a$ form the hyperplane $z(a)^\perp$. Hence $\xi$ lies in the kernel
exactly when $\xi\perp z(a)^\perp$, that is, when $\xi\in\R z(a)$. For the
translated form, let $S$ be the invertible matrix with
$Sz(X)=(1,(d^\alpha)_{1\le|\alpha|\le2})$. Then $Sz(a)$ is the first unit
vector, so the symmetric matrix $S^{-\mathsf T}\mathcal T_a(A)S^{-1}$ has
first row and column zero and a positive definite complementary block.

\emph{(c)} Since $A\succeq\mu I$, the matrix $\bar A=A-\mu\diag(I_n,0)$
satisfies $\bar A\succeq\mu\diag(0,I_{n^2})\succeq0$. The map
$B\mapsto\mathcal T_a(B)$ is linear, and \eqref{eq:heights-taylor-form} with
$B=\diag(I_n,0)$ gives
$z^{\mathsf T}\mathcal T_a(\diag(I_n,0))z=\int_0^1(1-t)\norm{d}^2dt=\frac12
\norm d^2$, because the first block of $U+tV$ is $d$. Therefore part (a)
gives
\[
 f=z^{\mathsf T}\mathcal T_a(\bar A)z+f(a)+b^{\mathsf T}d
   +\tfrac\mu2\norm d^2
  =z^{\mathsf T}\mathcal T_a(\bar A)z+\tfrac\mu2\norm{d+b/\mu}^2+c_a.
\]
As $z_0=1$, this is $f=z^{\mathsf T}\mathcal Q_az$. All three summands of
$\mathcal Q_a$ are positive semidefinite when $c_a\ge0$. Suppose $c_a>0$. The
first $n$ rows of $C_V(a)$ vanish, so
$\mathcal T_a(\bar A)\succeq\frac1{36}C_V(a)^{\mathsf T}\bar AC_V(a)
\succeq\frac\mu{36}C_V(a)^{\mathsf T}C_V(a)$, and
\[
 \mathcal Q_a\succeq\tfrac\mu{36}C_V(a)^{\mathsf T}C_V(a)
  +\tfrac\mu2C_{\rm aff}(a)^{\mathsf T}C_{\rm aff}(a)+c_ae_0e_0^{\mathsf T}.
\]
A vector $\xi$ in the kernel of the right side is orthogonal to the
coefficient vectors of every $d_id_j$, every $d_i+b_i/\mu$, and the constant
polynomial $1$. These polynomials span all polynomials of degree at most two,
because $d_i=(d_i+b_i/\mu)-(b_i/\mu)\cdot1$ and the constant, the $d_i$, and
the $d_id_j$ form a basis. Hence $\xi=0$ and $\mathcal Q_a\succ0$.

\emph{(d)} The entries of $C_U(a)$ and $C_V(a)$ are integer polynomials in
the coordinates of $a$. The vector $b=\nabla f(a)$ has entries in $E$, so
$C_{\rm aff}(a)$ and $c_a$ have entries in $E$ when $\mu\in E$. The matrices
$\mathcal T_a(A)$ and $\mathcal Q_a$ are obtained from these and from $A$ by
ring operations and rational constants. At a critical point $b=0$, and part
(a) reads $f-f(a)=z^{\mathsf T}\mathcal T_a(A)z$.

\emph{(e)} Substituting $A=\sum_ks_k\ell_k\ell_k^{\mathsf T}$ into the integral
of part (a) gives
$\sum_ks_k\int_0^1(1-t)(u_k+tv_k)^2dt$ with $u_k=\ell_k^{\mathsf T}U$ and
$v_k=\ell_k^{\mathsf T}V$. The scalar identity
$\int_0^1(1-t)(u+tv)^2dt=\frac12(u+v/3)^2+(v/6)^2$ and
$\frac12y^2=2(y/2)^2$ give the displayed sum. If $A\succeq0$ is rational,
rational symmetric elimination gives
$A=\sum_ks_k\ell_k\ell_k^{\mathsf T}$ with rational $s_k\ge0$ and rational
$\ell_k$, and each positive $s_k$ is a sum of rational squares
(Lemma~\ref{lem:models-rational-squares}): if
$s_k=\sum_l\sigma_{kl}^2$ with rational $\sigma_{kl}$, then
$s_ky^2=\sum_l(\sigma_{kl}y)^2$, and $2y^2=y^2+y^2$. Every resulting square
factor is a rational multiple of $\ell_k^{\mathsf T}(U+V/3)$ or of
$\ell_k^{\mathsf T}V$. The entries of $U$ and $V$ are $d_i$, $a_id_j$, and
$d_id_j$, whose coefficients are integer polynomials of degree at most two in
the coordinates of $a$, and $\ell_k$ is rational; this proves the statement
about coefficients. At a critical point $b=0$, so $f-f(a)$ is a sum of squares
of polynomials with coefficients in $\Q(a_1,\ldots,a_n)$.
\end{proof}

\subsection{The realization interface and the cube-root chain}
\label{app:heights-cube}

All quartics in Section~\ref{sec:heights} are assembled by
Lemma~\ref{lem:quartic-realization}, applied with as many residuals as
variables. We restate its hypotheses in the form used below; our lower
curvature bound $m$ is the parameter denoted $\mu$ in that lemma. Let $n\ge1$
and $p\in\R^n$, and let $G,r_1,\ldots,r_n$ be rational quadratic polynomials
vanishing at $p$, with exact translated forms
\[
 G(p+u)=\ell^{\mathsf T}u+u^{\mathsf T}Hu,\qquad
 r_j(p+u)=b_j^{\mathsf T}u+u^{\mathsf T}T_ju .
\]
Suppose positive rationals $m,\Lambda,\beta,\nu$ and a rational $t>0$ satisfy
\begin{gather}
 mI\preceq H,\quad \norm H\le\Lambda,\quad \norm{T_j}\le1,\quad
 \norm{b_j}\le\beta,\quad \sum_jb_jb_j^{\mathsf T}\succeq\nu^2I,\quad
 \norm\ell\le\varepsilon:=t^2,
 \label{eq:heights-R1}\\
 \varepsilon\le\min\Bigl\{1,\frac{m^2}{2n},
   \frac{\nu^2m^2}{36n(\Lambda+n\beta)^2}\Bigr\}.
 \label{eq:heights-R2}
\end{gather}
Then $F=(G/(t\nu))^2+\sum_{j=1}^n(r_j/\nu)^2$ is a rational quartic of
degree exactly four, a sum of $n+1$ squares of rational quadratics, with
$\nabla^2F\succeq\frac32I$ on $\R^n$; the point $p$ is its unique zero and
unique minimizer; and the canonical Hessian Gram of the displayed squares is a
rational positive definite full Hessian Gram of $F$, computed from the
rational factors by polynomially many rational operations on numbers of the
input bit length. No further approximation of $p$ is needed for the Gram.

\begin{lemma}[The cube-root chain]
\label{lem:heights-cube-chain}
Let $k\ge1$, $n=2k$, $M_k=1000^{k+3}$, and let $\delta_i$, $\xi_i$, and $p$ be
as in \eqref{eq:heights-tiny-chain} and \eqref{eq:heights-cube-gates}.
Then $0<\delta_i\le\delta_{i-1}^2$ for $1\le i\le k$; in particular
$0<\delta_k\le M_k^{-2^k}$ and $1<\xi_i\le1+M_k^{-2}$. The number $\xi_1$ is
irrational. In time polynomial in $k$ one can construct rational quadratics
$F_{k,1},\ldots,F_{k,n+1}$ and a rational matrix $A_k$, all of binary length
polynomial in $k$, such that $F_k=\sum_jF_{k,j}^2$ satisfies
$\nabla^2F_k\succeq\frac32I$, $F_k(p)=0$, and $F_k(X)>0$ for $X\ne p$, and
$A_k$ is a positive definite full Hessian Gram of $F_k$.
\end{lemma}

\begin{proof}
\emph{The chain.} Since $(1+\delta_i)^3=1+3\delta_{i-1}^2>1$, every
$\delta_i$ is positive. As $(1+t)^3\ge1+3t$ for $t\ge0$,
$1+3\delta_i\le1+3\delta_{i-1}^2$. Hence
$\delta_i\le\delta_0^{2^i}=M_k^{-2^i}$, which is at most
$M_k^{-2}\le10^{-24}$ for $i\ge1$. The gate equations
\eqref{eq:heights-cube-gates} follow from
$1+3\delta_{i-1}^2=1+3(\xi_{i-1}-1)^2$. For irrationality, note that
$M_k^2=K_0^3$ with $K_0=10^{2(k+3)}$, so $\xi_1=(K_0^3+3)^{1/3}/K_0$. A rational
cube root of an integer is an integer, and $K_0^3<K_0^3+3<(K_0+1)^3$, so
$\xi_1$ is irrational.

\emph{Residuals.} Put $c_1=1+3M_k^{-2}$ and
$c_i(X)=3y_{i-1}-6x_{i-1}+4$ for $2\le i\le k$, and define the $n$ rational
quadratics
\[
 r_i=x_i^2-y_i,\qquad s_i=x_iy_i-c_i\qquad(1\le i\le k).
\]
On their common real zero set $y_i=x_i^2$ and $x_i^3=c_i$. By induction on
$i$ and uniqueness of real cube roots, $x_i=\xi_i$; hence $p$ is the unique
common real zero. With $u=X-p$ and coordinates $u_{x_i},u_{y_i}$,
\begin{align*}
 r_i(p+u)&=2\xi_iu_{x_i}-u_{y_i}+u_{x_i}^2,\\
 s_i(p+u)&=\xi_i^2u_{x_i}+\xi_iu_{y_i}+[i\ge2]\,(6u_{x_{i-1}}-3u_{y_{i-1}})
          +u_{x_i}u_{y_i}.
\end{align*}
The quadratic parts are $e_{x_i}e_{x_i}^{\mathsf T}$ and
$\frac12(e_{x_i}e_{y_i}^{\mathsf T}+e_{y_i}e_{x_i}^{\mathsf T})$, of norms one and
one half. Since $1<\xi_i<1.001$, the gradients at $p$ satisfy
$\norm{\nabla r_i(p)}^2=4\xi_i^2+1<6$ and
$\norm{\nabla s_i(p)}^2\le\xi_i^4+\xi_i^2+45<48$, so $\beta=7$ bounds all of
them. The Jacobian $J$, with rows
$\nabla r_1(p)^{\mathsf T},\nabla s_1(p)^{\mathsf T},\ldots$ and columns ordered
as $X$, is block lower triangular with diagonal blocks
$\left(\begin{smallmatrix}2\xi_i&-1\\\xi_i^2&\xi_i\end{smallmatrix}\right)$ of
determinant $3\xi_i^2>3$. Hence $|\det J|>1$, and
$\norm J^2\le\norm J_F^2<54k\le(8k)^2$. Since $|\det J|$ is the product of the
singular values, $\sigma_{\min}(J)\ge|\det J|/\norm J^{n-1}>(8k)^{-(n-1)}$.
Thus $\sum(\nabla r_i\nabla r_i^{\mathsf T}+\nabla s_i\nabla s_i^{\mathsf T})
=J^{\mathsf T}J\succeq\nu^2I$ with $\nu=(8k)^{-(2k-1)}$.

\emph{An exact exposing quadratic.} For $1\le i\le k$ let
\[
 E_i=y_i^2-c_ix_i-\xi_is_i+\xi_i^2r_i .
\]
Each of $y_i^2-c_ix_i$, $s_i$, and $r_i$ vanishes at $p$, hence so does $E_i$.
Expanding,
$E_i=y_i^2-\xi_ix_iy_i+\xi_i^2x_i^2-\xi_i^2y_i-c_i(x_i-\xi_i)$. For a real
$\alpha$ let
$P_\alpha(x,y)=\alpha^2(x-\alpha)^2+\alpha(x-\alpha)(y-\alpha x)+(y-\alpha x)^2$,
which expands to $y^2-\alpha xy+\alpha^2x^2-\alpha^2y-\alpha^3x+\alpha^4$.
Therefore
\[
 E_i=P_{\xi_i}(x_i,y_i)+(\xi_i-x_i)(c_i-\xi_i^3).
\]
At $p+u$ we have $x_i-\xi_i=u_{x_i}$,
$y_i-\xi_ix_i=u_{y_i}-\xi_iu_{x_i}$, and
$c_i(p+u)-\xi_i^3=c_i(p+u)-c_i(p)$, which is $-6u_{x_{i-1}}+3u_{y_{i-1}}$ for
$i\ge2$ and zero for $i=1$. Hence
\begin{equation}
 E_i(p+u)=\xi_i^2u_{x_i}^2-\xi_iu_{x_i}u_{y_i}+u_{y_i}^2
   +[i\ge2]\,(6u_{x_i}u_{x_{i-1}}-3u_{x_i}u_{y_{i-1}}).
 \label{eq:heights-cube-exposer}
\end{equation}
In particular $E_i$ has neither a constant nor a linear term at $p$. The
local matrix $H_i=\left(\begin{smallmatrix}\xi_i^2&-\xi_i/2\\-\xi_i/2&1
\end{smallmatrix}\right)$ satisfies $\frac12I\preceq H_i\preceq2I$: the matrix
$H_i-\frac12I$ has positive diagonal and determinant $(\xi_i^2-1)/4\ge0$, and
$\tr H_i<2.01$.

Let $\omega_i=2^{-10(i-1)}$ and $E_*=\sum_i\omega_iE_i$, so that
$E_*(p+u)=u^{\mathsf T}H_*u$. Write $U_i=(u_{x_i},u_{y_i})$ and
$\tilde U_i=\sqrt{\omega_i}\,U_i$. The local terms contribute
$\sum_i\tilde U_i^{\mathsf T}H_i\tilde U_i$, which lies between
$\frac12\norm{\tilde U}^2$ and $2\norm{\tilde U}^2$. The cross term of index
$i\ge2$ has absolute value at most
\[
 \sqrt{45}\,\omega_i|U_i||U_{i-1}|
 =\sqrt{45}\,2^{-5}|\tilde U_i||\tilde U_{i-1}|
 \le\tfrac{\sqrt{45}}{64}\bigl(|\tilde U_i|^2+|\tilde U_{i-1}|^2\bigr),
\]
since $\omega_i/\omega_{i-1}=2^{-10}$. Summing over $i$, all cross terms
together have absolute value at most
$\frac{\sqrt{45}}{32}\norm{\tilde U}^2<0.21\norm{\tilde U}^2$. Hence
$\frac14\norm{\tilde U}^2\le u^{\mathsf T}H_*u\le\frac52\norm{\tilde U}^2$,
and $\omega_k\norm u^2\le\norm{\tilde U}^2\le\norm u^2$ gives
\begin{equation}
 \tfrac{\omega_k}4I\preceq H_*\preceq\tfrac52I.
 \label{eq:heights-cube-margin}
\end{equation}

\emph{Rational coefficients.} Let $0<\eta\le10^{-3}$, and let $\hat\xi_i$ be
rationals with $|\hat\xi_i-\xi_i|\le\eta$. Define the rational quadratic
\[
 \hat G=\sum_{i=1}^k\omega_i\bigl(y_i^2-c_ix_i-\hat\xi_is_i
          +\hat\xi_i^2r_i\bigr).
\]
The difference
$\hat G-E_*=\sum_i\omega_i(-(\hat\xi_i-\xi_i)s_i+(\hat\xi_i^2-\xi_i^2)r_i)$
is a combination of residuals, so $\hat G(p)=0$ exactly. Write
$\hat G(p+u)=\ell^{\mathsf T}u+u^{\mathsf T}Hu$. Since
$|\hat\xi_i^2-\xi_i^2|\le\eta(2\xi_i+\eta)\le3\eta$ and
$\sum_i\omega_i<1.001$,
\[
 \norm\ell\le\sum_i\omega_i\bigl(\eta\norm{\nabla s_i(p)}
     +3\eta\norm{\nabla r_i(p)}\bigr)<15\eta.
\]
The matrix $H-H_*$ is block diagonal, and its block $i$ is
$\omega_i(-(\hat\xi_i-\xi_i)T_{s_i}+(\hat\xi_i^2-\xi_i^2)T_{r_i})$, where
$T_{r_i},T_{s_i}$ are the quadratic parts above. Hence
$\norm{H-H_*}\le\frac72\eta$.

Put $m=\omega_k/8$, $\Lambda=3$, $\beta=7$, and $\nu=(8k)^{-(2k-1)}$. Let $t=2^{-\tau}$ for the least integer $\tau\ge0$ such that
$\varepsilon=4^{-\tau}$ satisfies \eqref{eq:heights-R2}, and let $\eta$ be the
largest power of two with $\eta\le\min\{10^{-3},\omega_k/32,\varepsilon/15\}$.
Then \eqref{eq:heights-cube-margin} and the perturbation bounds give
$H\succeq(\frac14-\frac7{64})\omega_kI\succeq mI$,
$H\preceq(\frac52+\frac7{64})I\preceq\Lambda I$, and
$\norm\ell<15\eta\le\varepsilon$. Hence \eqref{eq:heights-R1} and
\eqref{eq:heights-R2} hold for $\hat G$ and the residuals
$r_1,s_1,\ldots,r_k,s_k$.

\emph{Computing $\hat\xi_i$.} Let $\phi(s)=(1+3s^2)^{1/3}-1$, so that
$\delta_i=\phi(\delta_{i-1})$ and $|\phi'(s)|=2|s|(1+3s^2)^{-2/3}\le2|s|$. Put
$\hat\delta_0=\delta_0$. Given $\hat\delta_{i-1}$, compute by bisection on
$[1,2]$, with exact rational comparisons of $\rho^3$ against
$1+3\hat\delta_{i-1}^2$, a dyadic $\rho_i$ within $\eta/2$ of
$(1+3\hat\delta_{i-1}^2)^{1/3}$, and put $\hat\delta_i=\rho_i-1$. If
$|\hat\delta_{i-1}-\delta_{i-1}|\le\eta$, then
$|\hat\delta_{i-1}|,|\delta_{i-1}|\le M_k^{-1}+\eta<\frac1{500}$, so the mean
value theorem gives
\[
 |\hat\delta_i-\delta_i|\le\tfrac\eta2
   +|\phi(\hat\delta_{i-1})-\phi(\delta_{i-1})|
 \le\tfrac\eta2+\tfrac2{500}\eta\le\eta.
\]
By induction $\hat\xi_i=1+\hat\delta_i$ satisfies $|\hat\xi_i-\xi_i|\le\eta$
for all $i$. The numbers $\log_2(1/\nu)$, $\log_2(1/m)$, $\log_2(1/\varepsilon)$,
and $\log_2(1/\eta)$ are $O(k\log k)$. Every $\hat\delta_i$ is a dyadic rational
with $O(k\log k)$ bits, each bisection uses $O(k\log k)$ comparisons of
rationals of that length, and so all $\hat\xi_i$ are computed in time
polynomial in $k$.

\emph{Assembly.} The realization interface yields
$F_k=(\hat G/(t\nu))^2+\sum_i\bigl((r_i/\nu)^2+(s_i/\nu)^2\bigr)$, a sum of
$n+1$ squares with $\nabla^2F_k\succeq\frac32I$ and unique zero $p$, and the
rational positive definite canonical Hessian Gram $A_k$ of these factors.
Each factor has $O(k)$ monomials with coefficients of $O(k\log k)$ bits, so
$F_k$, its square factors, and $A_k$ have binary length polynomial in $k$.
Since $F_k$ is a sum of squares that vanish simultaneously only where all
residuals vanish, $F_k(X)>0$ for $X\ne p$.
\end{proof}

\subsection{Proof of Theorem~\ref{thm:heights-witness}}
\label{app:heights-witness}

\begin{proof}[Proof of Theorem~\ref{thm:heights-witness}]
Lemma~\ref{lem:models-hessian-gram}(a) and $B_k\succeq I$ give
$\nabla^2g_k\succeq I$. Since $F_k\ge0$ vanishes at $p$, we have
$\nabla F_k(p)=0$, so
\[
 g_k(p)=-\delta_k^2<0,\qquad \nabla g_k(p)=-2\delta_k\,e_{x_k},
 \qquad \norm{\nabla g_k(p)}=2\delta_k.
\]
For any $X$ with $r=\norm{X-p}$, strong convexity gives
$g_k(X)\ge-\delta_k^2-2\delta_kr+\frac12r^2$. If $g_k(X)\le0$, the quadratic
inequality forces $r\le(2+\sqrt6)\delta_k$. Hence every point of $S_k$,
including every boundary point, satisfies
\begin{equation}
 \norm{X-p}\le(2+\sqrt6)\,\delta_k<5M_k^{-2^k}.
 \label{eq:heights-localization}
\end{equation}
Thus $S_k$ is bounded; it is closed and convex, and it contains the open set
$\{g_k<0\}\ni p$. For distinct $X,Y\in S_k$ and $0<t<1$, strong convexity
gives $g_k(tX+(1-t)Y)<0$, so $S_k$ is strictly convex. The coordinates of
$p$ lie in $(1,1+3M_k^{-2}]$, and $5M_k^{-2^k}\le5\cdot10^{-24}$, which proves
the box inclusion.

For (b), let $X\in S_k$ have rational first coordinate $a/b$. The first
coordinate of $p$ is $\xi_1$, which is irrational
(Lemma~\ref{lem:heights-cube-chain}; $M_k^2$ is a perfect cube and
$M_k^2+3$ is not). Hence $(a/b)^3\ne\xi_1^3=1+3M_k^{-2}$, and
$M_k^2a^3-(M_k^2+3)b^3$ is a nonzero integer. Both $a/b$ and $\xi_1$ lie in
$(0.99,1.01)$, so \eqref{eq:heights-localization} gives
\begin{align*}
 1\le\bigl|M_k^2a^3-(M_k^2+3)b^3\bigr|
 &=M_k^2b^3\Bigl|\frac ab-\xi_1\Bigr|
   \Bigl|\frac{a^2}{b^2}+\frac ab\,\xi_1+\xi_1^2\Bigr|\\
 &<M_k^2b^3\cdot5M_k^{-2^k}\cdot3.1
 <16\,b^3M_k^{2-2^k}.
\end{align*}
Taking logarithms gives the displayed bound, since
$\log_2M_k=(k+3)\log_21000$. For $k\ge2$ we have $2^k-2\ge2^{k-1}$ and
$\log_21000>9.96$, so the right side exceeds
$2^{k-1}(3.3k+9.9)-2>k2^k$. Part (c) is Proposition~\ref{prop:upper-circuit-witness}, applied to
$g_k$ with the curvature bound $\mu=1$ certified by $B_k\succeq I$, since
$\min g_k\le g_k(p)<0$. The size of $(g_k,B_k)$ is polynomial in $k$ by
Lemma~\ref{lem:heights-cube-chain} and Lemma~\ref{lem:models-det-trace}.
\end{proof}

\subsection{Circle chains and the proof of Theorem~\ref{thm:rational-height}}
\label{app:heights-circle}

Both quartics of Theorem~\ref{thm:rational-height} come from one family of
circle chains. Fix $k\ge0$, $n=2(k+1)$, and a ratio $\rho\in\{1,\frac12\}$.
Put $R_j=\rho^j$ and $c_j=R_j/R_{j-1}^2$ for $j\ge1$; thus $c_j=1$ when
$\rho=1$ and $c_j=2^{j-2}$ when $\rho=\frac12$. Let
$z_j=((3+4\mathrm i)/5)^{2^j}=a_j+\mathrm ib_j$, and define the rational
residuals
\begin{equation}
 r_0=x_0-\tfrac35,\quad s_0=y_0-\tfrac45,\quad
 r_j=x_j-c_j(x_{j-1}^2-y_{j-1}^2),\quad s_j=y_j-2c_jx_{j-1}y_{j-1}
 \quad(1\le j\le k),
 \label{eq:heights-scaled-residuals}
\end{equation}
which reduce to \eqref{eq:heights-circle-residuals} when $\rho=1$. Put
$P=(R_ja_j,R_jb_j)_{j=0}^k$, $q_j=x_j^2+y_j^2-R_j^2$,
$\kappa_j=2c_j^2R_{j-1}^2$, and $t_{j-1}=(-b_{j-1},a_{j-1})$. Write
$u=X-P$, with $u_j\in\R^2$ the displacement of the pair $(x_j,y_j)$.

\begin{lemma}[Circle chains]
\label{lem:heights-circle-chain}
In this setting the following hold.
\begin{enumerate}[label=(\alph*)]
\item $P$ is the unique common real zero of the residuals, and every $q_j$
  vanishes at $P$.
\item For $1\le j\le k$ the rational-coefficient quadratic
  $E_j=q_j-2R_ja_j\,r_j-2R_jb_j\,s_j-\kappa_jq_{j-1}$ satisfies
  \[
   E_j(P+u)=\norm{u_j}^2-2\kappa_j\,(t_{j-1}^{\mathsf T}u_{j-1})^2 .
  \]
  Here $\kappa_j=2$ when $\rho=1$ and $\kappa_j=\frac12$ when $\rho=\frac12$.
\item The Jacobian of $(r_0,s_0,\ldots,r_k,s_k)$ at $P$ is $J=I-T$, where $T$
  is block strictly lower triangular, its only nonzero blocks are
  $T_{j,j-1}=2c_jR_{j-1}\left(\begin{smallmatrix}a_{j-1}&-b_{j-1}\\
  b_{j-1}&a_{j-1}\end{smallmatrix}\right)$, and $2c_jR_{j-1}$ equals $2$ when
  $\rho=1$ and $1$ when $\rho=\frac12$. In particular $T^{k+1}=0$.
\end{enumerate}
\end{lemma}

\begin{proof}
\emph{(a)} The residuals force $(x_0,y_0)=(3/5,4/5)$ and then, inductively,
$x_j+\mathrm iy_j=c_j(x_{j-1}+\mathrm iy_{j-1})^2$. Since
$c_j(R_{j-1}z_{j-1})^2=c_jR_{j-1}^2z_j=R_jz_j$, the unique solution is $P$.
Also $|R_jz_j|=R_j$, so $q_j(P)=0$.

\emph{(b)} Fix $j\ge1$ and write $(a,b)=(a_{j-1},b_{j-1})$,
$(\alpha',\beta')=R_{j-1}(a,b)$ for the predecessor point,
$(A',B')=R_j(a_j,b_j)=c_j(\alpha'^2-\beta'^2,2\alpha'\beta')$ for the successor
point, $c=c_j$, $u_{j-1}=(u_1,u_2)$, and $u_j=(w_1,w_2)$. Then
\begin{align*}
 q_j(P+u)&=2A'w_1+2B'w_2+\norm{u_j}^2,\\
 r_j(P+u)&=w_1-c(2\alpha'u_1-2\beta'u_2+u_1^2-u_2^2),\\
 s_j(P+u)&=w_2-2c(\beta'u_1+\alpha'u_2+u_1u_2),\\
 q_{j-1}(P+u)&=2\alpha'u_1+2\beta'u_2+\norm{u_{j-1}}^2 .
\end{align*}
In $E_j=q_j-2A'r_j-2B's_j-\kappa_jq_{j-1}$ the terms linear in $u_j$ cancel,
and there are no products of $u_j$ with $u_{j-1}$. The identities
$A'\alpha'+B'\beta'=c\,\alpha'R_{j-1}^2$ and
$-A'\beta'+B'\alpha'=c\,\beta'R_{j-1}^2$ show that the terms linear in
$u_{j-1}$ are $(4c^2R_{j-1}^2-2\kappa_j)(\alpha'u_1+\beta'u_2)=0$. The
quadratic terms in $u_{j-1}$ have matrix
$2c\left(\begin{smallmatrix}A'&B'\\B'&-A'\end{smallmatrix}\right)-\kappa_jI$.
Since $2cR_j=2c^2R_{j-1}^2=\kappa_j$, this matrix equals
$\kappa_j\left(\begin{smallmatrix}a_j-1&b_j\\b_j&-a_j-1\end{smallmatrix}\right)$,
and $a_j=a^2-b^2$, $b_j=2ab$, $a^2+b^2=1$ turn it into
$-2\kappa_j\left(\begin{smallmatrix}b^2&-ab\\-ab&a^2\end{smallmatrix}\right)
=-2\kappa_jt_{j-1}t_{j-1}^{\mathsf T}$. The constant term vanishes by (a). The
values of $\kappa_j$ follow from $c_j=1$, $R_{j-1}=1$ when $\rho=1$, and from
$c_jR_{j-1}=2^{j-2}\cdot2^{-(j-1)}=\frac12$ when $\rho=\frac12$.

\emph{(c)} The derivatives of $(r_j,s_j)$ with respect to $(x_{j-1},y_{j-1})$
at $P$ form the matrix
$-2c_j\left(\begin{smallmatrix}\alpha'&-\beta'\\\beta'&\alpha'\end{smallmatrix}
\right)$, which is $-T_{j,j-1}$; the derivatives with respect to $(x_j,y_j)$
form the identity, and no other derivative is nonzero.
\end{proof}

\begin{proof}[Proof of Theorem~\ref{thm:rational-height}]
Both quartics are assembled by the realization interface of
Appendix~\ref{app:heights-cube}.

\emph{Approximating the chain.} Let $0<\eta\le1$ be rational. Put
$\hat z_0=z_0$ and let $\hat z_j$ be the complex number obtained from
$\hat z_{j-1}^2$ by rounding its real and imaginary parts to multiples of
$h=2^{-P'}$, where $P'$ is the least integer with $3^kh\le\eta$. Let
$e_j=|\hat z_j-z_j|$. If $e_{j-1}\le1$, then
$|\hat z_{j-1}^2-z_{j-1}^2|\le(|\hat z_{j-1}|+|z_{j-1}|)e_{j-1}\le3e_{j-1}$,
and rounding adds at most $\sqrt2h<2h$, so $e_j\le3e_{j-1}+2h$. Since
$e_0=0$, induction gives $e_j\le(3^j-1)h\le\eta$ for all $j\le k$. For
$j\ge1$ the real and imaginary parts of $\hat z_j$ are dyadic rationals of
absolute value at most two with $O(k+\log(1/\eta))$ bits; the initial parts
$\frac35$ and $\frac45$ of $\hat z_0$ have constant bit length. The
computation takes time polynomial in $k+\log(1/\eta)$. Multiplying by $R_j$
gives rational approximations of $R_ja_j$ and $R_jb_j$ with errors at most
$\eta$.

\emph{The exposing quadratic.} With weights $\omega_j>0$ to be chosen, let
$E_0=r_0^2+s_0^2$ and
\[
 \hat G=\omega_0E_0+\sum_{j=1}^k\omega_j\bigl(q_j-2\hat\alpha_jr_j
   -2\hat\beta_js_j-\kappa_jq_{j-1}\bigr),
\]
where $\hat\alpha_j,\hat\beta_j$ are the approximations of $R_ja_j,R_jb_j$.
Only coefficients of residuals are approximated, so $\hat G(P)=0$. Let $H_*$
be the matrix of the exact sum $\omega_0E_0+\sum_j\omega_jE_j$; by
Lemma~\ref{lem:heights-circle-chain}(b) it is block diagonal, with block
$\omega_jI-2\kappa_{j+1}\omega_{j+1}t_jt_j^{\mathsf T}$ for $j<k$ and
$\omega_kI$ for $j=k$. Write $\hat G(P+u)=\ell^{\mathsf T}u+u^{\mathsf T}Hu$,
and let $\Delta_j=(\hat\alpha_j-R_ja_j,\hat\beta_j-R_jb_j)$, so
$\norm{\Delta_j}\le\sqrt2\eta$. The quadratic parts of $r_j$ and $s_j$ are
$-c_j\diag(1,-1)$ and $-c_j\left(\begin{smallmatrix}0&1\\1&0\end{smallmatrix}
\right)$ on the pair $j-1$. Hence $H-H_*$ is block diagonal with blocks of
norm $2\omega_jc_j\norm{\Delta_j}\le2\sqrt2\,\omega_jc_j\eta$, and
$\norm\ell\le\sum_{j\ge1}2\omega_j\norm{\Delta_j}_1
\max\{\norm{\nabla r_j(P)},\norm{\nabla s_j(P)}\}$. By
Lemma~\ref{lem:heights-circle-chain}(c),
$\norm{\nabla r_j(P)}=\norm{\nabla s_j(P)}=(1+(2c_jR_{j-1})^2)^{1/2}$ for
$j\ge1$, and $\norm{\nabla r_0}=\norm{\nabla s_0}=1$.

\emph{The quartic $f_k$ ($\rho=1$).} Take $\omega_j=8^{-j}$. Each block of
$H_*$ with $j<k$ is at least $(\omega_j-4\omega_{j+1})I=\frac{\omega_j}2I$
and at most $\omega_jI$, so $8^{-k}I\preceq H_*\preceq I$. The residual
gradients have norms at most $\sqrt5<3$, and
$\norm{H-H_*}\le2\sqrt2\eta/8<\eta/2$,
$\norm\ell\le4\sqrt5\,\eta\sum_{j\ge1}8^{-j}<2\eta$. By
Lemma~\ref{lem:heights-circle-chain}(c), $\norm T\le2$ and
$J^{-1}=\sum_{i=0}^kT^i$, so $\norm{J^{-1}}<2^{k+1}$ and
$J^{\mathsf T}J\succeq4^{-(k+1)}I$. Put $m=8^{-k}/2$, $\Lambda=2$, $\beta=3$, and
$\nu=2^{-(k+1)}$. Choose the dyadic $t$ as in the
cube-root construction, so that $\varepsilon=t^2$ satisfies
\eqref{eq:heights-R2}, and take $\eta=\min\{8^{-k},\varepsilon/2\}$. Then
$H\succeq mI$, $H\preceq\Lambda I$, and $\norm\ell\le\varepsilon$, so the
realization interface applies to $\hat G$ and the $n$ residuals. It yields
\[
 f_k=\Bigl(\frac{\hat G}{t\nu}\Bigr)^2
     +\sum_{j=0}^k\Bigl[\Bigl(\frac{r_j}\nu\Bigr)^2+\Bigl(\frac{s_j}\nu\Bigr)^2\Bigr],
\]
a sum of $n+1$ rational squares with $\nabla^2f_k\succeq\frac32I$, unique
zero and minimizer $p=P$, and a rational positive definite Hessian Gram.
Here $\log_2(1/m)$, $\log_2(1/\nu)$, $\log_2(1/\varepsilon)$, and
$\log_2(1/\eta)$ are $O(k)$, so every coefficient has $O(k)$ bits and the
construction takes time polynomial in $k$. Part (b) follows from
Lemma~\ref{lem:heights-circle-denominators} with $m=2^k$, because the
last two coordinates of $p$ are $a_k$ and $b_k$.

\emph{The quartic $\tilde f_k$ ($\rho=\frac12$).} Take $\omega_j=k+1-j$. Each
block of $H_*$ with $j<k$ is $\omega_jI-\omega_{j+1}t_jt_j^{\mathsf T}$, whose
eigenvalues are $\omega_j$ and $\omega_j-\omega_{j+1}=1$, and the last block
is $I$; hence $I\preceq H_*\preceq(k+1)I$. Let $C=\max\{1,2^{k-2}\}$, so
$c_j\le C$. The residual gradients have norm $\sqrt2$ for $j\ge1$, hence
$\norm{H-H_*}\le2\sqrt2\,kC\eta<3kC\eta$ and
$\norm\ell\le4\sqrt2\,\eta\sum_{j\ge1}\omega_j<3k(k+1)\eta$. Divide every
residual by $C$: the quadratic parts of $r_j/C$ and $s_j/C$ have norm
$c_j/C\le1$, their gradients have norm at most $\beta=2/C$, and, since
$\norm T\le1$ by Lemma~\ref{lem:heights-circle-chain}(c),
$\norm{J^{-1}}\le k+1$, so $J^{\mathsf T}J/C^2\succeq\nu^2I$ with
$\nu=1/(C(k+1))$. Put $m=\frac12$ and $\Lambda=k+2$. Choose $t$ as before
and $\eta=\min\{1,1/(6kC),\varepsilon/(3k(k+1))\}$ (for $k=0$ no
approximation is needed). Then $H\succeq\frac12I$, $H\preceq(k+2)I$, and
$\norm\ell\le\varepsilon$, and the interface yields
\[
 \tilde f_k=\Bigl(\frac{\hat G}{t\nu}\Bigr)^2
   +\sum_{j=0}^k\Bigl[\Bigl(\frac{r_j}{C\nu}\Bigr)^2
     +\Bigl(\frac{s_j}{C\nu}\Bigr)^2\Bigr]
 =\Bigl(\frac{\hat G}{t\nu}\Bigr)^2+(k+1)^2\sum_{j=0}^k(r_j^2+s_j^2),
\]
with $\nabla^2\tilde f_k\succeq\frac32I$, unique zero and minimizer
$\tilde p=P$, and a rational positive definite Hessian Gram. All
coefficients have $O(k\log k)$ bits. The coordinates of $\tilde p$ are
$2^{-j}a_j$ and $2^{-j}b_j$; dividing by a power of two cannot remove the
factor $5^{2^k}$ from the reduced denominators of $a_k$ and $b_k$, and
$\norm{\tilde p}^2=\sum_j4^{-j}<\frac43$.

\emph{The Hessian at $\tilde p$.} All square factors vanish at $\tilde p$, so
$\nabla^2(\phi^2)(\tilde p)=2\nabla\phi(\tilde p)\nabla\phi(\tilde p)^{\mathsf T}$
for each factor $\phi$, and $C\nu=1/(k+1)$ gives
\[
 \nabla^2\tilde f_k(\tilde p)=\frac{2\ell\ell^{\mathsf T}}{\varepsilon\nu^2}
   +2(k+1)^2J^{\mathsf T}J .
\]
Since $\norm T\le1$, we have $\norm J\le2$ and $\norm{J^{-1}}\le k+1$, so the
second term lies between $2I$ and $8(k+1)^2I$. The first term is positive
semidefinite with norm at most
$2\norm\ell^2/(\varepsilon\nu^2)\le2\varepsilon/\nu^2
\le m^2/(18n(\Lambda+n\beta)^2)<1$ by \eqref{eq:heights-R2}. This proves (c).
Finally, every ordinary representation of $p$ or $\tilde p$ contains a
coordinate whose reduced denominator is at least $5^{2^k}$, and
$\log_25^{2^k}=2^k\log_25$; the inputs have length polynomial in $k$.
\end{proof}

\subsection{Moments, optimal Grams, and exposing matrices}
\label{app:heights-moment}

We use two elementary facts about positive semidefinite matrices. If
$P,R\succeq0$ and $\tr(PR)=0$, then
$\tr(PR)=\norm{R^{1/2}P^{1/2}}_F^2$ forces $R^{1/2}P^{1/2}=0$, hence $PR=0$
and $\operatorname{range}R\subseteq\ker P$. If $P\succeq0$ and
$v^{\mathsf T}Pv=0$, then $Pv=0$.

\begin{proof}[Proof of Theorem~\ref{thm:heights-moment}]
\emph{(a)} Since $\nabla f(p)=0$, Lemma~\ref{lem:taylor-sos} with $a=p$ gives
$f-f_*=z^{\mathsf T}Q_*z$, where $Q_*=\mathcal T_p(A)\succeq0$,
$\ker Q_*=\R e$ by part (b) of that lemma, and the entries of $Q_*$ lie in
$K$ by part (d), because $f$ and $A$ are rational. If $Q$ is any positive
semidefinite Gram of $f-f_*$, then $e^{\mathsf T}Qe=f(p)-f_*=0$, so $Qe=0$ and
$\rank Q\le N-1$.

\emph{(b)} For $y$ feasible in (P) and $(\gamma,Q)$ feasible in (S),
\[
 \Lambda_y(f)-\gamma=\Lambda_y(z^{\mathsf T}Qz)
 =\sum_{\alpha,\beta}Q_{\alpha\beta}y_{\alpha+\beta}=\tr(QM_2(y))\ge0 .
\]
The point sequence $y^p_\alpha=p^\alpha$ is feasible in (P), with
$M_2(y^p)=ee^{\mathsf T}$ and value $f(p)=f_*$, and $(f_*,Q_*)$ is feasible in
(S). Hence both optimal values equal $f_*$ and both are attained. If $y$ is
optimal in (P), then $\tr(Q_*M_2(y))=\Lambda_y(f)-f_*=0$, so the range of
$M_2(y)$ lies in $\ker Q_*=\R e$, and $M_2(y)=c\,ee^{\mathsf T}$ for some
$c\ge0$. The constant entry gives $c=y_0=1$. Every multi-index $\gamma$ with
$|\gamma|\le4$ is a sum $\alpha+\beta$ with $|\alpha|,|\beta|\le2$, so
$y_\gamma=M_2(y)_{\alpha\beta}=p^\gamma$, and $y=y^p$. The ranks of
$ee^{\mathsf T}$ and $Q_*$ add up to $N$, so the pair is strictly
complementary. The moments of the standard Gaussian measure on $\R^n$ are
strictly feasible in (P): for $c\ne0$,
$c^{\mathsf T}M_2(y)c$ is the Gaussian integral of the square of the nonzero
polynomial $c^{\mathsf T}z$, which is positive. In (S), the pair
$(f_*-1,Q_*+e_0e_0^{\mathsf T})$ is feasible, and $Q_*+e_0e_0^{\mathsf T}\succ0$:
a nonzero $\xi\notin\R e$ has $\xi^{\mathsf T}Q_*\xi>0$, while $\xi=ce$ with
$c\ne0$ has $(e_0^{\mathsf T}\xi)^2=c^2>0$ because the constant coordinate of
$e$ is one.

\emph{(c)} If $K\subseteq E$, then $Q_*$ is an optimal Gram of rank $N-1$ with
entries in $E$. Conversely, let $Q$ be an optimal Gram of rank $N-1$ with
entries in $E$. Its rank is the same over $E$ and over $\R$, so its kernel in
$E^N$ is spanned by one nonzero vector $\zeta$. By (a), $\zeta\in\R e$, and
since the constant coordinate of $e$ is one, $\zeta_0\ne0$ and
$e=\zeta/\zeta_0\in E^N$. The coordinates of $e$ include $p_1,\ldots,p_n$,
so $K\subseteq E$.

\emph{(d)} Let $V=\Span_\Q\{p^\alpha:|\alpha|\le2\}\subset\R$, of dimension
$r$ over $\Q$. The $\Q$-linear map $\Q^N\to V$, $\xi\mapsto\xi^{\mathsf T}e$, is
onto, so its kernel has dimension $N-r$. If $Q$ is a rational optimal Gram,
then $Qe=0$ by (a), so every row of $Q$ lies in this kernel, and the rank of
$Q$, which is its row rank over $\Q$, is at most $N-r$. If some $p_i$ is
irrational, then $1$ and $p_i$ are linearly independent over $\Q$, so $r\ge2$
and $\rank Q+\rank(ee^{\mathsf T})\le N-1$.

\emph{(e)} The set $\mathcal F_p=\{Q\succeq0:Qe=0\}$ is the face of the
positive semidefinite cone exposed by $ee^{\mathsf T}$, and
$\mathcal G_*\subseteq\mathcal F_p$ by (a). The minimal face of the positive
semidefinite cone containing a matrix $Q_0\succeq0$ consists of all
$Q\succeq0$ with $\operatorname{range}Q\subseteq\operatorname{range}Q_0$. For
$Q_0=Q_*$ the range is $e^\perp$, so this minimal face is $\mathcal F_p$, and
$\mathcal F_p$ is the smallest face containing $\mathcal G_*$; $Q_*$ lies in
its relative interior. If $Z\succeq0$ exposes $\mathcal G_*$, then
$\tr(ZQ_*)=0$, so $\operatorname{range}Z\subseteq\R e$ and
$Z=c\,ee^{\mathsf T}$ with $c>0$ because $Z\ne0$. Conversely
$\tr(ee^{\mathsf T}Q)=e^{\mathsf T}Qe=0$ for every $Q\in\mathcal G_*$. For
$Z=c\,ee^{\mathsf T}$, the entry in the row of the constant monomial and the
column of $X_i$ is $cp_i$, and the constant entry is $c$, so
$p_i=Z_{0,X_i}/Z_{0,0}$ belongs to the field generated by the entries of $Z$.
In particular no nonzero rational positive semidefinite matrix exposes
$\mathcal G_*$ when $p\notin\Q^n$.

For facial reduction, let $\mathcal A$ map a symmetric matrix $Q$ to the
coefficient vector of $z^{\mathsf T}Qz$, indexed by the monomials of degree at
most four. The optimal-level Gram system is
$\mathcal A(Q)=\operatorname{coef}(f-f_*)$, $Q\succeq0$. For every vector $y$
indexed by those monomials,
$\ip y{\mathcal A(Q)}=\sum_{\alpha,\beta}Q_{\alpha\beta}y_{\alpha+\beta}
=\ip{M_2(y)}Q$, so $\mathcal A^*(y)=M_2(y)$. For $y=y^p$ we obtain
$\mathcal A^*(y^p)=ee^{\mathsf T}$, which is positive semidefinite and nonzero,
and $\ip{y^p}{\operatorname{coef}(f-f_*)}=f(p)-f_*=0$. This is a
facial-reduction certificate in the sense of Borwein and Wolkowicz
\cite{BorweinWolkowicz1981}: every feasible $Q$ satisfies
$\ip{ee^{\mathsf T}}Q=\ip{y^p}{\mathcal A(Q)}=0$, so the system may be
restricted to $\mathcal F_p$. After this one step the feasible point $Q_*$
lies in the relative interior of $\mathcal F_p$, so no further step is
needed. Zero steps do not suffice, because no feasible Gram is positive
definite. Hence the real singularity degree, the least number of such steps
after which the restricted system has a relatively interior feasible point,
is exactly one.
\end{proof}

\begin{proof}[Proof of Corollary~\ref{cor:heights-gram}]
Let $p$ be the optimizer of $f_k$. By Theorem~\ref{thm:rational-height}(b)
and Lemma~\ref{lem:heights-circle-denominators}, its last two coordinates are
$a_k$ and $b_k$, with reduced denominators exactly $5^{2^k}$.

\emph{(a)} By Theorem~\ref{thm:heights-moment}(b) the optimal moment matrix
is $ee^{\mathsf T}$, whose entry in the row of the constant monomial and the
column of $x_k$ is $a_k$.

\emph{(b)} The Taylor Gram $Q_*$ of Theorem~\ref{thm:heights-moment}(a) is
rational, since $p$ and the Hessian Gram are rational, and has rank $N-1$.
Let $Q$ be any rational optimal Gram of rank $N-1$. Order $z$ with the
constant monomial first and write
\[
 Q=\begin{pmatrix}q_{00}&c^{\mathsf T}\\c&Q'\end{pmatrix},
 \qquad e=\begin{pmatrix}1\\w\end{pmatrix}.
\]
The block $Q'$ is positive definite: for $u\ne0$ the vector $(0,u)$ is not a
multiple of $e$, so $(0,u)^{\mathsf T}Q(0,u)=u^{\mathsf T}Q'u>0$ by
Theorem~\ref{thm:heights-moment}(a). The equation $Qe=0$ gives $Q'w=-c$. For
each row $i$ of the $(N-1)\times N$ matrix $[Q'\mid c]$, multiply by the
product $\pi_i$ of the reduced denominators of its $N$ entries. This produces
an integer system $Q''w=-c''$ with $Q''=\diag(\pi)Q'$ nonsingular. By
Cramer's rule every coordinate of $w$ is a ratio of an integer to
$\det Q''$, so its reduced denominator divides $|\det Q''|$. One coordinate
of $w$ is $a_k$, hence
\begin{equation}
 5^{2^k}\le|\det Q''| .
 \label{eq:heights-cramer}
\end{equation}
If every entry of $Q$ has height at most $B$, then $\pi_i\le2^{NB}$, every
entry of $Q''$ has absolute value at most $2^{(N+1)B}$, and the Leibniz
expansion gives $|\det Q''|\le(N-1)!\,2^{(N+1)(N-1)B}$. With
\eqref{eq:heights-cramer} this is the stated bound on $B$. For the total
length, let $s_i$ be the sum, over the entries $a/b$ of row $i$ of
$[Q'\mid c]$, of $\log_2b$, plus the largest $\log_2\max\{1,|a|\}$ in that
row. The entries of row $i$ of $Q''$ have absolute value at most $2^{s_i}$, so
$|\det Q''|\le(N-1)!\,2^{\sum_is_i}$. Since $s_i$ is at most the sum of the
binary lengths of the entries in row $i$, \eqref{eq:heights-cramer} shows that
the total binary length of the entries of $Q$ is at least
$H_k-\log_2((N-1)!)$. As $\log_2((N-1)!)=O(n^2\log n)$ and
$N^2-1=\Theta(n^4)$, the entry bound is $\Omega(2^k/n^4)$.

\emph{(c)} By Theorem~\ref{thm:heights-moment}(e), a nonzero positive
semidefinite exposing matrix is $Z=c\,ee^{\mathsf T}$ with $c>0$, and
$a_k=Z_{0,x_k}/Z_{0,0}$. If these two entries are $\alpha_1/\beta_1$ and
$\alpha_2/\beta_2$ in lowest terms with heights at most $B$, then
$a_k=\alpha_1\beta_2/(\beta_1\alpha_2)$, whose reduced denominator divides
$\beta_1\alpha_2$ and is therefore at most $2^{2B}$. Hence
$5^{2^k}\le2^{2B}$.

\emph{(d)} If $f_k=\sum_{j=1}^{n+1}F_j^2$ with rational quadratics $F_j$, then
$\sum_j\operatorname{coef}(F_j)\operatorname{coef}(F_j)^{\mathsf T}$ is a
rational positive semidefinite Gram of $f_k=f_k-0$ of rank at most $n+1$ and
binary length polynomial in $k$. Since $n\ge2$, $n+1<N-1$.

For $\tilde f_k$ the same proofs apply, because the terminal coordinates of
$\tilde p$ have reduced denominators divisible by $5^{2^k}$: in (b) and (c)
these denominators divide $|\det Q''|$ and $\beta_1\alpha_2$, respectively, so
both quantities are still at least $5^{2^k}$.
\end{proof}

\subsection{Interior Gram matrices}
\label{app:heights-interior}

\begin{lemma}[A uniform trace bound]
\label{lem:heights-cube-moments}
Let $W_n=2^{-n}\int_{[-1,1]^n}z(X)z(X)^{\mathsf T}dX$. Then
$W_n\succeq I_N/(15(n+1))$. Consequently every positive semidefinite Gram
$Q$ of a polynomial $h$ of degree at most four satisfies
$\tr Q\le15(n+1)\norm h_1$.
\end{lemma}

\begin{proof}
For $\xi\in\R^N$, $\xi^{\mathsf T}W_n\xi$ is the mean of $P^2$ over the cube,
where $P=\xi^{\mathsf T}z=a_0+\sum_ib_iX_i+\sum_ic_iX_i^2
+\sum_{i<j}d_{ij}X_iX_j$. Under the uniform probability measure on
$[-1,1]^n$, the functions $1$, $X_i$, $X_i^2-\frac13$, and $X_iX_j$ ($i<j$)
are pairwise orthogonal, with mean squares $1$, $\frac13$, $\frac4{45}$, and
$\frac19$. Writing $a'=a_0+\frac13\sum_ic_i$,
\[
 \xi^{\mathsf T}W_n\xi=a'^2+\tfrac13\sum_ib_i^2+\tfrac4{45}\sum_ic_i^2
   +\tfrac19\sum_{i<j}d_{ij}^2 .
\]
By the Cauchy--Schwarz inequality,
$a_0^2\le2a'^2+\frac29(\sum_ic_i)^2\le2a'^2+\frac{2n}9\sum_ic_i^2$, so
\[
 \norm\xi^2\le2a'^2+\sum_ib_i^2+\bigl(1+\tfrac{2n}9\bigr)\sum_ic_i^2
   +\sum_{i<j}d_{ij}^2 .
\]
Each coefficient on the right is at most $15(n+1)$ times the corresponding
coefficient of $\xi^{\mathsf T}W_n\xi$; for the squares this is
$1+\frac{2n}9\le\frac43(n+1)$. This proves $W_n\succeq I/(15(n+1))$. If
$Q\succeq0$ and $h=z^{\mathsf T}Qz$, then
$\tr Q/(15(n+1))\le\tr(QW_n)$, which is the mean of $h$ over the cube and
hence at most $\max_{[-1,1]^n}|h|\le\norm h_1$.
\end{proof}

\begin{lemma}[Denominators of a positive determinant]
\label{lem:heights-det-denominators}
Let $Q$ be a rational symmetric matrix whose upper-triangular entries have
reduced denominators $b_{ij}$, $i\le j$. Then
$\bigl(\prod_{i\le j}b_{ij}^2\bigr)\det Q$ is an integer. If $Q\succ0$, then
$\det Q\ge\prod_{i\le j}b_{ij}^{-2}$.
\end{lemma}

\begin{proof}
In the Leibniz expansion of $\det Q$, a product
$\prod_iQ_{i,\sigma(i)}$ uses each diagonal entry at most once and each
off-diagonal entry $Q_{ij}=Q_{ji}$ at most twice, as $Q_{ij}$ and as $Q_{ji}$.
Hence every term, and therefore $\det Q$, becomes an integer after
multiplication by $\prod_{i\le j}b_{ij}^2$. For $Q\succ0$ that integer is
positive, hence at least one.
\end{proof}

\begin{proof}[Proof of Theorem~\ref{thm:interior-gram}]
\emph{(b)} By Lemma~\ref{lem:heights-cube-chain}, $A_k\succ0$ is rational, so
$\mu_{A_k}>0$, $\lambda_kA_k\succeq3I$, and
$\lambda_kA_k+2\iota_k\iota_k^{\mathsf T}\succeq3I$; this matrix is a full Hessian Gram of
$h_k$ because the Hessian biform of $u_k^2$ is $2v_{x_k}^2$. Hence
$\nabla^2h_k\succeq3I$ (Lemma~\ref{lem:models-hessian-gram}(a)). The binary
length of $\lambda_k$ is polynomial in $k$ by
Lemma~\ref{lem:models-det-trace}, so the whole certified quartic has length
polynomial in $k$. Since $F_k\ge0$ vanishes only at $p$ and
$u_k(p)=\delta_k>0$, $h_k>0$ on $\R^n$, and its minimum $m_k$ is attained by
coercivity, so $m_k>0$. As $\nabla h_k(p)=2\delta_ke_{x_k}\ne0$, the point
$p$ is not the minimizer, and
$m_k<h_k(p)=\delta_k^2\le M_k^{-2^{k+1}}$. If $F_{k,j}$ are the square factors
of $F_k$, the matrix
$\lambda_k\sum_j\operatorname{coef}(F_{k,j})\operatorname{coef}(F_{k,j})^{\mathsf T}
+\operatorname{coef}(u_k)\operatorname{coef}(u_k)^{\mathsf T}$ is a rational
positive semidefinite Gram of $h_k$ of rank at most $n+2$, and $n+2<N$ for
$n\ge2$. Writing $\lambda_k=\sum_{\epsilon\in S}2^\epsilon$ in binary and
$2^\epsilon F^2=(2^{\epsilon/2}F)^2$ or
$2^\epsilon F^2=2(2^{(\epsilon-1)/2}F)^2$ turns
$\lambda_k\sum_jF_{k,j}^2+u_k^2$ into an unweighted sum of $O((n+1)\log_2\lambda_k)$
squares of rational quadratics. Finally, apply
Lemma~\ref{lem:taylor-sos}(c) to $h_k$ with $\mu=3$: the number
$c_a=h_k(a)-\norm{\nabla h_k(a)}^2/6$ depends continuously on $a$ and equals
$m_k>0$ at the minimizer of $h_k$, so it is positive at some rational $a$, and
$\mathcal Q_a$ is a rational positive definite Gram of $h_k$.

\emph{(c)} Let $Q\succ0$ be a rational Gram of $h_k$. Since the constant
coordinate of $z(p)$ is one and $z(p)$ has further nonzero coordinates,
$\norm{z(p)}>1$ and
\[
 \lambda_{\min}(Q)\le\frac{z(p)^{\mathsf T}Qz(p)}{\norm{z(p)}^2}<h_k(p)
   \le M_k^{-2^{k+1}}.
\]
By Lemma~\ref{lem:heights-cube-moments}, every eigenvalue of $Q$ is at most
$\tr Q\le T_k$, and $T_k\ge1$. Therefore
$0<\det Q<M_k^{-2^{k+1}}T_k^{N-1}$. Lemma~\ref{lem:heights-det-denominators}
gives $\det Q\ge\prod_{i\le j}b_{ij}^{-2}$, and taking logarithms proves the
displayed inequality. The polynomial $h_k$ has polynomially many coefficients
of polynomial binary length in $k$, so $\log_2T_k$ is polynomial in $k$, and
$N=O(k^2)$. Since $\log_2M_k=(k+3)\log_21000$, the right side is
$\Omega(k2^k)$. There are $N(N+1)/2=O(k^4)$ upper-triangular entries, so one
of them has $\Omega(2^k/k^3)$ denominator bits.

\emph{(a)} Put $\mu=\mu_A$, which has binary length polynomial in $L$
(Lemma~\ref{lem:models-det-trace}), and
$\varphi=f-\norm{\nabla f}^2/(2\mu)$. This polynomial has degree at most six,
$O(n^6)$ coefficients, and coefficients of binary length polynomial in $L$.
At the minimizer $p$ of $f$, $\varphi(p)=f(p)=\min f>0$. Let $\Delta$ be the
product of the denominators of the coefficients of $\varphi$; then
$\Delta\varphi$ has integer coefficients of binary length $\tau=\poly(L)$, and
$\{\Delta\varphi>0\}=\{\varphi>0\}$ is a nonempty open set. By the theorem of
Basu, Pollack, and Roy on rational points in sets defined by strict
inequalities \cite[Theorem~4.1.2]{BasuPollackRoy1996}, it contains a point
$a\in\Q^n$ whose coordinates have numerators and denominators of binary length
at most $\tau\,6^{O(n)}=\poly(L)\,2^{O(n)}$. Then $c_a=\varphi(a)>0$, and
$\mathcal Q_a$ of Lemma~\ref{lem:taylor-sos}(c) is a rational positive
definite Gram of $f$.

Each entry of $\mathcal Q_a$ is the value at $a$ of a polynomial of degree at
most six with $\poly(n)$ terms and rational coefficients of binary length
$\poly(L)$; these polynomials do not depend on $a$, and they come from $f$,
$A$, and $\mu$ through the fixed formulas \eqref{eq:heights-taylor-gram} and
\eqref{eq:heights-completed-gram}. If the coordinates of $a$ have binary
length at most $B_a$, write them over a common denominator of length at most
$nB_a$; then every such value has binary length $\poly(L)\cdot B_a$. Summing
over the $N^2$ entries gives total length $\poly(L)\,2^{O(n)}$. A bound that
only counted arithmetic operations would not suffice here; the bound uses the
fixed degree of these expressions in $a$.

For the unweighted sum of squares, rational $LDL^{\mathsf T}$ elimination of
$\mathcal Q_a\succ0$ gives $f=\sum_i d_i(\ell_i^{\mathsf T}z)^2$ with rational
$d_i>0$, and the entries of $L$ and the $d_i$ are ratios of minors of
$\mathcal Q_a$, hence of binary length polynomial in that of $\mathcal Q_a$.
Each $d_i$ is a sum of at most $2\bits(d_i)$ rational squares
(Lemma~\ref{lem:models-rational-squares}), which gives the claimed
expression. The argument proves existence and a size bound. It is not an
algorithm polynomial in $L$, because the sampling theorem is applied to a set
whose rational points may need exponentially long coordinates.
\end{proof}

\subsection{Shared-circuit Grams}
\label{app:heights-circuit}

\begin{proof}[Proof of Theorem~\ref{thm:circuit-gram}]
\emph{Validation of the input.} Check that $f$ has degree at most four and
that $A$ is symmetric, verify \eqref{eq:heights-hessian-gram} by expanding
$w^{\mathsf T}Aw$ and comparing coefficients with
$v^{\mathsf T}\nabla^2f(X)v$, and verify $A\succ0$ by rational
$LDL^{\mathsf T}$ elimination. These are polynomial-time rational
computations. Reject on failure.

\emph{The observable and the Newton point.} Compute $\mu=\mu_A$, which has
binary length polynomial in the input length
(Lemma~\ref{lem:models-det-trace}) and satisfies $\nabla^2f\succeq\mu I$
(Lemma~\ref{lem:models-hessian-gram}(a)), and compute the expanded polynomial
$h=f-\norm{\nabla f}^2/(2\mu)$. It has degree at most six, $O(n^6)$
coefficients, and binary length polynomial in the input length. Since
$\nabla f(p)=0$, we have $h(p)=f(p)=\min f$. Apply
Lemma~\ref{lem:newton-circuit}(a) to $f$ and $\mu$, with $L'$ at least the
encoding lengths of the input of that lemma and of $h$. It gives, in
polynomial time and without oracle calls, a shared rational circuit for a
point $\hat x\in\Q^n$, in which every division is by a nonzero value, and a
number $g=g_{L'}>0$ such that $h(p)=0$ or $|h(p)|\ge g$, and
$|h(\hat x)-h(p)|\le g/8$.

\emph{The Gram.} Let $Q=\mathcal Q_{\hat x}$ be the matrix
\eqref{eq:heights-completed-gram} with center $\hat x$ and
$\bar A=A-\mu\diag(I_n,0)$. Its ingredients are circuits: the entries of
$C_U(\hat x)$ and $C_V(\hat x)$ are integer polynomials of degree at most two
in $\hat x$, $C_{\rm aff}(\hat x)$ uses $\nabla f(\hat x)/\mu$, and
$c_{\hat x}=f(\hat x)-\norm{\nabla f(\hat x)}^2/(2\mu)=h(\hat x)$. Evaluating
$\nabla f$ at $\hat x$ takes $O(n^4)$ gates, and the matrix products in
\eqref{eq:heights-taylor-gram} have dimensions $O(n^2)$ and take $O(n^6)$
gates. All are appended to the shared circuit for $\hat x$, and the only new
divisors are $\mu$ and the constants $2$, $3$, and $36$. By
Lemma~\ref{lem:taylor-sos}(c), $f=z^{\mathsf T}Qz$ for every center; no
approximation enters this identity, because $\hat x$ denotes an exact rational
vector.

If $\min f>0$, then $h(p)>0$, so $h(p)\ge g$ and
$c_{\hat x}=h(\hat x)\ge7g/8>0$, and Lemma~\ref{lem:taylor-sos}(c) gives
$Q\succ0$. Conversely, if $Q\succ0$, then for every $X$,
$f(X)=z(X)^{\mathsf T}Qz(X)\ge\lambda_{\min}(Q)\norm{z(X)}^2\ge
\lambda_{\min}(Q)>0$, since the constant coordinate of $z(X)$ is one; as the
minimum is attained, $\min f>0$.

\emph{The remarks after the theorem.} For the family $h_k$ of
Theorem~\ref{thm:interior-gram}, $\min h_k>0$, so the value of $Q$ is a
rational positive definite Gram of $h_k$, whose expanded encoding has
$\Omega(k2^k)$ denominator bits. For item (i), the map from a certified
quartic to the circuit matrix $Q$ is a polynomial-time reduction from the
predicate $\min f>0$ to positive definiteness of circuit matrices, and that
predicate is $\PosSLP$-hard by Theorem~\ref{thm:quartic-complete}(i). For item
(ii), $f(\hat x)\ge\min f$ for every point, and
Proposition~\ref{prop:upper-circuit-witness} gives $f(\hat x)<0$ when
$\min f<0$. For item (iv), if $\min f>0$ the leading principal minors of $Q$
are positive, so symmetric elimination without pivoting divides only by
positive quantities and yields circuits for
$f=\sum_id_i(\ell_i^{\mathsf T}z)^2$ with $d_i>0$; when $\min f\le0$ some pivot
may vanish, so this step is not part of the sign-free construction.
\end{proof}
